%% ============================================================
%%  On Keeping Humans Meaningfully in the Loop of AI Agents:
%%  A Critical Review and Research Agenda
%%
%%  Author: Ahmet Selim Fedakar
%%  A response to and critical review of Mitchell, Ghosh & Passi,
%%  "AI Agents Push Humans Out of the Loop" (NeurIPS 2026, arXiv:2608.23642).
%%
%%  Build: pdflatex main && bibtex main && pdflatex main && pdflatex main
%% ============================================================
\documentclass[11pt]{article}

\usepackage[utf8]{inputenc}
\usepackage[T1]{fontenc}
\usepackage{lmodern}
\usepackage[margin=1in]{geometry}
\usepackage{microtype}
\usepackage{amsmath,amssymb}
\usepackage{graphicx}
\usepackage{booktabs}
\usepackage{enumitem}
\usepackage{xcolor}
\usepackage[numbers,sort&compress]{natbib}
\usepackage{titlesec}
\usepackage[colorlinks=true,linkcolor=black,citecolor=blue,urlcolor=blue]{hyperref}

% Editorial note macro: anything still to be sourced or decided.
\newcommand{\todo}[1]{\textcolor{red}{\textbf{[TODO: #1]}}}

\titleformat{\section}{\large\bfseries}{\thesection}{0.6em}{}
\titleformat{\subsection}{\normalsize\bfseries}{\thesubsection}{0.6em}{}

\title{\textbf{On Keeping Humans Meaningfully in the Loop of AI Agents:\\
A Critical Review and Research Agenda}}

\author{%
  Ahmet Selim Fedakar\\
  \small Independent Researcher\thanks{Affiliation and contact to be finalized before submission. \todo{set affiliation}}\\
  \small \texttt{a.selimfedakar@gmail.com}%
}

\date{\today}

\begin{document}
\maketitle

\begin{abstract}
\noindent
Autonomous AI agents are increasingly deployed with human oversight proposed as the
primary safeguard. In a recent NeurIPS position paper, \citet{mitchell2026agents} argue
that current approaches to building and deploying agent systems do not merely fail to
support human oversight---they contribute to its degradation: overseers face output
generated faster than they can review, approval fatigue sets in, and, invoking
\citeauthor{bainbridge1983}'s \emph{ironies of automation}, the very act of oversight
erodes the cognitive capacities that effective oversight requires. This paper (i)
reconstructs that argument and situates it in the human--automation interaction
literature, from \citeyear{bainbridge1983} onward; (ii) critically evaluates how well its
central empirical claim---that extended AI use causes skill atrophy in knowledge work---is
supported, distinguishing a strongly grounded \emph{descriptive} lineage from a largely
\emph{asserted} applied claim; and (iii) proposes a concrete research agenda with testable
hypotheses for studying oversight degradation in agentic systems. We argue that the
position is directionally correct and well-motivated, but that its most consequential
claim remains an inference by analogy awaiting direct measurement---which we frame as the
field's central open problem.
\end{abstract}

% ============================================================
\section{Introduction}
\label{sec:intro}

``Keep a human in the loop'' has become the default governance answer to the risks of
increasingly autonomous AI agents. It appears in corporate deployment policy, in
regulatory text such as the EU AI Act's meaningful-control requirement, and in the design
of nearly every production agent framework, which gates consequential actions behind a
human approval step. The appeal is obvious: a competent human supervisor can, in
principle, catch an agent's mistakes before they cause harm.

A recent position paper by \citet{mitchell2026agents} challenges the premise directly. Its
claim is not the familiar one that oversight is \emph{hard}; it is the stronger and more
unsettling claim that the way we currently build and deploy agents \emph{actively degrades}
the human capacities oversight depends on. If that is right, then human-in-the-loop is not
merely an imperfect safeguard---it is a self-undermining one, strengthening the very
dependency it is meant to check.

This paper has three aims, corresponding to three things a researcher entering this area
needs. First (Section~\ref{sec:position}), a faithful reconstruction of the position's
argument and its proposed remedies. Second (Section~\ref{sec:landscape}), a map of the
intellectual lineage it stands on---four decades of human-factors research on automation
and a fast-moving 2025--2026 empirical frontier. Third (Sections~\ref{sec:critique}
and~\ref{sec:agenda}), a critical evaluation that separates what is well-established from
what is asserted, and a research agenda built on that gap. Throughout, we are deliberately
careful about evidence level: a central contribution of this review is to show \emph{which}
parts of the argument are load-bearing empirical fact and which are, as yet, well-motivated
conjecture.

% ============================================================
\section{The Position: Mitchell, Ghosh \& Passi}
\label{sec:position}

\subsection{Provenance and framing}

The target is ``AI Agents Push Humans Out of the Loop'' \citep{mitchell2026agents}, a
position paper accepted to NeurIPS (arXiv:2608.23642; v1 24 August 2026, v3 6 September
2026; cs.AI/cs.HC), authored by Margaret Mitchell and Avijit Ghosh of Hugging Face and
Samir Passi of the Data \& Society Research Institute. It is explicitly a \emph{position}
paper: it advances a normative argument and a design stance rather than reporting a new
experiment. It is accompanied by a Data \& Society primer, ``The Oversight Fallacy: Why AI
Agents Require More than Humans-in-the-Loop'' \citep{passi2026oversight}, which grounds a
closely related argument in fieldwork in a computational biology laboratory and enumerates
four components of oversight in practice---knowledge, observation, control, and
intervention.

\subsection{The core argument}

The position has two linked halves, stated plainly in the abstract. The first is that
current agent \emph{design} impedes effective oversight. The second---and the paper's
distinctive contribution---is that the cognitive capacities oversight requires are
\emph{themselves degraded} by extended use of AI systems. Together these yield the thesis
that current development and deployment approaches ``do not support effective human
oversight---they contribute to its degradation.'' The paper frames part of this degradation
as \emph{passively incentivized} by prevailing design and deployment choices rather than
actively intended, a nuance worth preserving when restating the claim.

The mechanism the paper emphasizes is concrete. Traditional safeguards such as approval
prompts work only when a user can recognize a mistake and intervene before its consequences
cascade. Agents routinely violate that precondition: they generate voluminous output faster
than a human can meaningfully review it, and the result is that overseers find themselves
``repeatedly approving prompts,'' producing what the paper calls \emph{approval fatigue}.

\subsection{The irony of automation: ``oversight degrades the overseer''}

The paper devotes a full section (its Section~4, ``The irony of automation'') to grounding
this in \citet{bainbridge1983}. Bainbridge's classic observation is that the more capable an
automated system becomes, the more the human operator's skill and situational awareness
atrophy through disuse---so that operators are least prepared to intervene precisely when
intervention matters most. Mitchell, Ghosh, and Passi extend this from the control room to
the AI overseer and compress it into a memorable formulation: ``the very act of being an
overseer degrades the capacities oversight requires: \emph{oversight degrades the
overseer}.'' This is the conceptual heart of the paper.

\subsection{The prescriptive core: cognitive scaffolding}

The paper's remedy is a normative demand plus a two-stage solution set it calls
\emph{cognitive scaffolding}, ``provided by developers and deployers.'' The normative demand
is that the human needs of overseers be treated as equal in importance to AI-agent
capability---a reweighting of priorities, not merely an add-on. The two stages are:

\begin{itemize}[leftmargin=*,itemsep=2pt]
  \item \textbf{Development-stage design affordances:} \emph{strategic friction}
        (deliberately slowing the human at decision points), \emph{decision design}, and
        \emph{behavioral monitoring}.
  \item \textbf{Deployment-stage organizational protocols:} trainings and exercises, role
        design, audits, and attention to workload and scheduling.
\end{itemize}

These are organized against a set of oversight \emph{goals} in a goals-versus-mechanisms
table: the paper identifies six high-level goals---preventing automation bias, fatigue,
overreliance, and cognitive overload, and supporting judgement and engagement---and uses
their intersection with solution categories to catalog existing mechanisms and locate gaps.
We note two simplifications to avoid overstating the structure: decision design spans both
stages, and the deployment side carries a workload/scheduling category, so the clean
three-and-three split is a convenience rather than the paper's exact taxonomy.

% ============================================================
\section{The Research Landscape}
\label{sec:landscape}

The position's power comes from standing on a deep and well-established human-factors
tradition. This section maps that lineage and the recent frontier the paper joins.

\subsection{The ironies of automation}

\citet{bainbridge1983} is the keystone. Its core ``irony'' is double: automating the easy
parts of a task leaves the human the hard residual, and the skills needed to handle that
residual decay precisely because the automation now handles the routine. The result is an
operator who is both less practiced and more essential---a structural trap, not an
individual failing. The paper's ``oversight degrades the overseer'' is a direct descendant.

\subsection{Levels of automation and the human's changing role}

Bainbridge's insight was formalized into the \emph{levels of automation} framework.
\citet{sheridan1978} introduced a ten-level scale from fully manual to fully autonomous
control; \citet{parasuraman2000} generalized this into a model of \emph{types and levels}
of human interaction with automation across four stages (information acquisition, analysis,
decision, action). This taxonomy---one of the most cited in the field---gives a vocabulary
for exactly where a human sits relative to an agent, and therefore for what ``in the loop''
even means in a given system.

\subsection{Automation bias, complacency, and the out-of-the-loop problem}

Two robust empirical phenomena underpin the degradation concern.
\citet{parasuraman1997}, in ``Humans and Automation: Use, Misuse, Disuse, Abuse''---one of
the most-cited papers in human factors---established how people systematically mis-calibrate
their trust in automation. \citet{parasuraman2010} then argued that \emph{automation
complacency} (insufficient monitoring) and \emph{automation bias} (over-accepting automated
cues) are not distinct phenomena but overlapping manifestations of the same effect, with
reallocation of attentional resources as the central shared mechanism. We flag one
precision point for anyone citing this: the 2010 paper is explicitly an \emph{integration},
and it names a second, non-attentional route (``active discounting,'' a decision bias), so
attentional withdrawal is the dominant but not sole cause. This body of work is the
strongest empirical anchor for the position's claims---and, as Section~\ref{sec:critique}
argues, also the source of its central vulnerability, because it was established on
rule-based automation in aviation and process control, not on generative AI agents.

\subsection{Deskilling and skill atrophy}

The step from ``operators lose skill through disuse'' (well-documented in aviation and
process control) to ``knowledge workers and developers lose skill through extended agent
use'' is the position's most consequential move. It is also, as we argue below, the least
directly evidenced in the agentic setting---the clearest opening for new research.

\subsection{Moral crumple zones and accountability}

\citet{elish2019} supplies the accountability counterpart to the cognitive one. A
\emph{moral crumple zone} is the human operator who absorbs blame for the failure of a
system they could not realistically control---structurally positioned to be held
responsible while lacking the capacity to intervene effectively. This sharpens why
degraded oversight is not only a safety problem but an injustice one: a human kept nominally
``in the loop'' primarily to bear liability is oversight as theater.

\subsection{Meaningful human control and the critique of mandated oversight}

\citet{santonidesio2018} articulate \emph{meaningful human control} as a philosophical
standard---tracking and tracing conditions a system must satisfy for human control to be
real rather than nominal. In parallel, a critical literature questions whether mandated
human oversight works at all: \citet{green2022} argues that policies requiring human
oversight of government algorithms often fail on their own terms, and a thread we group
under ``fictional human oversight'' \citep{fictionaloversight2024} contends that such
requirements can provide the appearance of control without the substance. The position
paper lands inside this live controversy, not a settled consensus.

\subsection{The 2025--2026 empirical frontier}

A cluster of very recent work makes the ``oversight is harder than frameworks assume'' case
empirically---and is where a new researcher should look for collaborators and open threads.
We stress an attribution point: these are \emph{distinct} papers from the position paper,
sharing at most one author, and should be cited to their own sources.

\begin{itemize}[leftmargin=*,itemsep=3pt]
  \item \textbf{Oversight in practice.} \citet{dhanorkar2026} interviewed 17 developers who
        use coding agents (e.g., Claude Code, Cline) at least weekly and found they perform
        \emph{at least four} forms of oversight---a priori control, co-planning, real-time
        monitoring, and post hoc review---showing oversight is not only reactive but
        preventative and proactive. Crucially, developers settle for ``efficient, not
        perfect'' oversight, exposing a gap between governance frameworks that assume
        oversight is feasible and the messy reality of using agents.
  \item \textbf{Verification can backfire.} \citet{recchia2025} show that when people did
        independent research to check an AI's answers, they became \emph{more} confident in
        those answers even when the answers were wrong---and this held even under
        substantial financial incentives for accuracy. Confirmation bias, in other words,
        can defeat the ``human verification scales'' assumption at the core of scalable
        oversight.
  \item \textbf{Oversight as good work.} \citet{faas2025} ground oversight in work-design
        psychology: effectiveness depends not only on error detection but on the overseer's
        motivation and the meaningfulness of the work. Through four co-design workshops they
        derive four requirements---understanding tasks and responsibilities, insight into
        the AI's decision-making, meaningful contribution, and collaboration with peers and
        the AI---integrated with the SMART work-design model.
  \item \textbf{When not to automate.} \citet{delachica2026} propose a pre-decision filter
        (PHP-AIO) that, they argue, acts \emph{before} the deployment-commitment choice,
        distinguishing itself from NIST, EU AI Act, and HITL taxonomies that engage only
        after a deployment candidate exists. We mark this medium-confidence: the paper's
        self-positioning is contestable (the NIST AI RMF MAP function and EU AI Act Art.~5
        do bear on whether to deploy), and its detailed internal structure could not be
        independently verified.
\end{itemize}

% ============================================================
\section{Critical Evaluation}
\label{sec:critique}

\subsection{What is well supported}

The position's \emph{descriptive} foundation is strong. Automation bias and complacency,
the out-of-the-loop performance problem, and the ironies of automation are established
across decades of peer-reviewed human-factors research
\citep{bainbridge1983,parasuraman1997,parasuraman2000,parasuraman2010}. That human
oversight of agents is harder than governance frameworks assume is now also empirically
supported in the agentic setting \citep{dhanorkar2026}, and that human verification can
actively backfire is directly demonstrated \citep{recchia2025}. On these points the paper
is not speculating; it is reporting a mature consensus and extending it with fresh evidence.

\subsection{The evidence gap: skill atrophy asserted, not measured}

The paper's most consequential claim---that extended use of AI agents specifically causes
\emph{skill atrophy} in knowledge and software work---is, in the agentic setting, largely
\emph{asserted by analogy} to the older automation literature rather than demonstrated by
direct longitudinal measurement. The companion empirical works establish adjacent facts:
that oversight is effortful and imperfect \citep{dhanorkar2026}, and that verification can
increase misplaced confidence \citep{recchia2025}. Neither establishes AI-induced
deskilling as a \emph{measured} effect over time. This is not a fatal objection to a
position paper---whose job is to argue, not to run the experiment---but it is the precise
boundary between what the field knows and what it suspects, and it should be stated as such.
Mistaking the analogy for a finding is exactly the overclaim this review is designed to
prevent.

\subsection{Does oversight add safety, or theater?}

If human verification can backfire \citep{recchia2025} and, in some reported settings,
intervention succeeds only a minority of the time, then the burden shifts: the interesting
question is not \emph{whether} to keep a human in the loop but \emph{under what design
conditions} doing so measurably adds safety. The position paper's proposed affordances
(strategic friction, decision design) are plausible answers, but they are hypotheses, not
yet validated interventions. Combined with the moral-crumple-zone critique
\citep{elish2019}, this is the strongest version of the skeptical case: oversight that
neither catches errors nor is designed to can still serve to allocate blame.

\subsection{What the paper does \emph{not} claim}

For accuracy, three things attributed to this literature in secondary summaries did
\emph{not} survive verification and should not be repeated as the position paper's content:
it does \emph{not} present a clean four-item list of named mechanisms (approval fatigue /
automation bias / loss of situational awareness / skill atrophy) as its taxonomy; it does
\emph{not} contain a ``twelve design considerations via SMART'' scheme (that framing belongs
to a different line of work, \citet{faas2025}, and even there the twelve-item operational
claim failed verification); and the PHP-AIO protocol's detailed gate/category structure
\citep{delachica2026} could not be verified and should not be stated as fact. Flagging these
guards the present paper against inheriting other summaries' errors.

% ============================================================
\section{A Research Agenda}
\label{sec:agenda}

The gap identified above is an opportunity. We propose four directions, each with a
testable core.

\subsection{Direction 1: Measure AI-induced skill atrophy directly}

\textbf{Open question.} Is skill atrophy in knowledge/agentic work a real, measurable
effect, or an inference from classic human-factors analogies?
\textbf{Testable hypothesis.} Over months of sustained agent use, a cohort's unaided
performance on the underlying task (e.g., debugging without the agent, or situational
awareness during a hand-back) declines relative to a matched control, controlling for task
difficulty. \textbf{Why it matters.} No direct longitudinal measurement is cited in the
literature; this is the field's central empirical hole. It is also safety-critical in
embodied settings such as a voice-AI automotive agent, where driver vigilance decay during
extended assisted driving is exactly the ``out-of-the-loop'' problem with lives attached.

\subsection{Direction 2: Identify design conditions under which oversight actually helps}

\textbf{Open question.} Given that verification can backfire \citep{recchia2025} and
intervention success can be low, when does oversight add safety rather than theater?
\textbf{Testable hypothesis.} The position paper's ``strategic friction'' and decision
design measurably raise intervention-success and error-catch rates in an agent workflow,
and do so without being neutralized by approval fatigue over time. \textbf{Method sketch.}
Controlled studies that manipulate friction and decision-surface design and measure
caught-error rate, false-approval rate, and their decay.

\subsection{Direction 3: Toward an integrated theory of oversight design}

\textbf{Open question.} How do the competing solution framings reconcile---the position
paper's ``cognitive scaffolding'' \citep{mitchell2026agents}, the work-design/SMART ``good
work'' framing \citep{faas2025}, and the ``sometimes don't automate'' pre-decision filter
\citep{delachica2026}? Which affordances are \emph{necessary} versus \emph{sufficient}?
\textbf{Why it matters.} This synthesis is currently unclaimed; a unifying design theory of
agent oversight is an opening a new researcher could own.

\subsection{Direction 4: Field, venue, and framing}

The active cluster spans Hugging Face (Mitchell, Ghosh), Data \& Society (Passi, Singh),
Microsoft Research (Dhanorkar, Vorvoreanu), and Saarland/Freiburg HCI (Langer, Feit, and
colleagues), building on the classic human-factors tradition (the Parasuraman lineage).
Relevant venues are the NeurIPS/ICML position tracks, FAccT, CSCW, AAAI, and \emph{Human
Factors}. \textbf{Strategic question for a new entrant.} Does the agentic-oversight gap
belong in ML safety, HCI, or human factors---and which framing gives the most leverage? Our
reading is that the measurement gap (Direction~1) is best pursued with HCI/human-factors
methods but published into the ML-safety conversation, where the claim currently circulates
without the measurement.

% ============================================================
\section{Conclusion}
\label{sec:conclusion}

\citet{mitchell2026agents} make a directionally correct and well-motivated argument: the
dominant ``human in the loop'' safeguard is in tension with how agents are actually built,
and the oversight relationship may erode the very capacities it relies on. The descriptive
core of that argument rests on decades of solid human-factors research and a credible new
empirical frontier. Its most striking claim---that extended AI-agent use causes skill
atrophy in knowledge work---remains, in the agentic setting, an inference awaiting direct
measurement. That gap is not a weakness to dismiss the paper over; it is the research
program the paper implicitly calls for. Keeping humans \emph{meaningfully} in the loop will
require treating overseer cognition as a first-class design target---and, first of all,
measuring whether the degradation we fear is actually happening.

% ============================================================
\section*{Acknowledgments}
\todo{Optional. Note the use of an automated literature-search and verification pipeline in
the methods/acknowledgments if the venue requires tool disclosure.}

\bibliographystyle{plainnat}
\bibliography{references}

\end{document}
